\section{ATrust and DCP-AI}\label{sec:atrust}

\ATrust is attributed to Edison Vazquez as an unpublished conceptual protocol
\citep{vazquez_atrust}.  In the conceptual model, an agent identity relates to
credential descriptions, handshake requirements, trust-ledger entries, and an
evidence map.  \Argorix implements syntax, semantic relationships, bytecode
representations, and trace events for these concepts.  It does not implement
real credential issuance or verification, cryptographic challenge--response,
remote DID resolution, or live attestation.

Identity declarations associate agents with identifier methods and
boundaries.  Credential contracts describe issuer/subject/type relationships.
Handshake declarations describe participants and required evidence.  A trust
ledger is a local, hash-linked declarative structure whose entries reference
subjects and evidence.  Evidence maps require named declarations to be covered
by artifact links.  Validation rejects dangling references and prohibited
claims.  The word \emph{handshake} therefore denotes a checked contract, not an
executed cryptographic exchange; \emph{ledger} denotes a local evidence data
structure, not blockchain consensus or immutable storage.

\DCPAI, attributed to Danilo Naranjo Emparanza, proposes a Digital Citizenship
Protocol for Autonomous AI and a control-plane framing for agents
\citep{naranjo_dcpai}.  It complements ATrust's trust relationships and
\Argorix's compiled enforcement boundary, but the three are not conflated.
DCP-AI is a published conceptual proposal; ATrust is unpublished conceptual
work; \Argorix is the software artifact evaluated here.

The relationship can be stated operationally.  DCP-AI motivates questions of
agent standing and governance; ATrust supplies a vocabulary for identity,
credentials, handshakes, ledgers, and evidence; \Argorix checks local
declarations and emits evidence about their handling.  None of these local
checks demonstrates that a remote party is genuine.

The evidence scope figure in \cref{fig:evidence-chain} distinguishes the
unverified source input from bundle-verified JSON artifacts.  Its scope does
not extend to an issuer, credential authority, or remote agent.

MCP and A2A bridges follow the same principle.  They bind agents, passports,
identities, message contracts, and boundaries while requiring network and
execution denial in the studied profile.  The declarations can support future
protocol adapters \citep{mcp_spec_2025,a2a_spec}; today they are declarative
boundaries and auditable negative controls.
